\chapter{Introduction}
\label{ch:introduction}



\section{Overview}
\label{sec:intro:overview}


Randomized control trials represent a fast-growing and critical tool in empirical economics, boasting a literature that is vast, still growing \citep{Imbens2023}, and has recently birthed its own 
Nobel prizes \citep{Banerjee2020,Duflo2020,Kremer2020}. Much of the focus on this literature, however, is on the internal validity of the experimental design, as opposed to the external validity of results for policy relevance \citep{Heckman2008,Deaton2018,Manski2013}. This thesis concerns itself with two key components of external validity: representativity and ecological validity \citep{Neyman1934,Schmuckler2001}. It develops novel methodologies for leveraging the particular properties of digital ad platforms to address both of these components: (i) ecological validity when studying the impact of digital interventions and (ii) representativity more broadly, through optimizing participant recruitment subject to budget contraints. 

\section{Methodologies for External validity}
\label{sec:intro:ev}


The relative paucity of focus on external validity has been noted by the field's own critics for decades. \citet{Manski2013} writes that ``analyses of experimental data have tended to be silent on the problem of extrapolating from the experiments performed to policies of interest'' (p.~37); the credibility of a study, in his account, depends jointly on internal and external validity (pp.~36--37). \citet{Heckman2008} echoes the diagnosis: forecasting ``the effects of interventions in environments other than the one where the estimate was obtained'' --- ``the problem of external validity'' (p.~8) --- is a task that randomised estimation, by construction, does not deliver. His structural account distinguishes between the statistical estimation problem (what RCTs solve) and the economic forecasting problem (what policy requires), and names the distance between them as the field's open question. \citet{Deaton2018} restates the critique in its strongest contemporary form: randomisation buys internal validity by design, but ``credibility in estimation can lead to incredibility in use'' (p.~3).

The recognition is not purely theoretical. A body of empirical work demonstrates that randomised estimates often fail to replicate across contexts, populations, and time periods (\cite{Pritchett2015, Allcott2015, Bisbee2017, Rosenzweig2019}). Additionally, the scale-up literature has documented the pilot-to-scale voltage drop empirically \citep{Bold2018} and programmatically \citep{Muralidharan2017,Al-Ubaydli2020}. Even from inside the randomised-controlled-trial movement itself, external validity is acknowledged as unfinished business. \citet{Snowberg2016} write that ``creating a rigorous framework for external validity is an important step in completing an ecosystem for social science field experiments.''

This thesis does not engage external validity in the abstract; it engages it for a specific and fast-growing class of interventions: digital interventions. Over the past two decades an increasing share of the contact between institutions and the public --- commercial, governmental, and humanitarian alike --- has migrated onto digital platforms: display and social-media advertising, mobile applications, mobile messaging, and search \citep{Goldfarb2019,Aker2010}. Public-interest actors have followed, repurposing the infrastructure built for commercial persuasion \citep{LewisRao2015} to deliver health information, behavioural nudges, and educational content at population scale \citep{Athey2023,Breza2021,Ho2023}. This migration is what makes digital interventions both the object of study and, as we will argue, the instrument of measurement throughout this thesis: the advertising platform that delivers a malaria-prevention message is also the platform on which a sample can be recruited; and the app store that distributes a parenting app is also a channel on which take-up and engagement become visible. Finally, this migration drives the need for us to approach these new platforms not just as interesting applications of existing theory, but to fruther our understanding of the theoretical goals in each stage of study preparation and delivery, and use those goals to measure the opportunities and challenges presented by these new environments in a rigorous manner. 



\section{Defining External Validity and Experiment Types}
\label{sec:intro:definitions}

\citet{CampbellStanley1963} decompose external validity across four sub-dimensions: populations, settings, treatment variables, and measurement variables. We will use the term \emph{representatvity} to address the dimension of ``population.'' A sample is representative of a population if its joint distribution on outcome-relevant covariates matches that of the target population --- the standard sense from the stratified-sampling tradition of \citet{Neyman1934}, where a covariate is considered relevant when it is correlated with the survey outcome (\citealt{Cochran1977}, pp.~99--100). It's critical to define what is meant by ``relevant'' covariates---notably this is not a generic or fixed concept but related to the question at hand. We will use the term \emph{ecological validity} to bundle the remaining three sub-dimensions from \citet{CampbellStanley1963}: the setting, stimulus, and response to the conditions of an intervention. This is aligned with the three-dimensional definition of the term in \citet{Schmuckler2001}.

\citet{Harrison2004} similarly define six properties to determine the field context of an experiment, which can broadly map to representativity and ecological validity: (i) the subject pool, (ii) the prior knowledge of the subject pool, (iii) the item, (iv) the task, (v) the stakes, (vi) the environment. They also lay out a taxonomy for field experiments which will be useful for defining what we will do in this thesis. They distinguish the \emph{conventional lab experiment}, conducted with a standard subject pool of students in an abstract setting, from three field variants of increasing ecological realism: the \emph{artefactual field experiment}, which retains the lab protocol but draws subjects from a non-standard population; the \emph{framed field experiment}, in which a non-standard subject pool engages with a real-world commodity, task, or information set, but does so in the knowledge that they are participating in a study; and the \emph{natural field experiment}, in which subjects undertake naturally occurring tasks in their naturally occurring environment, unaware that they are participants. The malaria feed experiment of Chapter~\ref{ch:malaria} is a natural field experiment on both the delivery and measurement sides --- treatment is a real ad in a real news feed, served by an entity with no apparent connection to the surveying organisation --- while the encouragement design of Chapter~\ref{ch:bebbo} is a framed field experiment in the Harrison--List sense: real subjects engage with a real app under real-world adoption conditions, but they have consented to and are aware of the study.

\section{Digital Recruitment}
\label{sec:intro:digital-recruitment}

Across the experimental social sciences --- including behavioural and experimental economics, psychology, and political science --- the default subject pool has shifted away from university student samples and, in survey research, away from telephone-based probability sampling, toward internet-based crowdsourcing platforms such as Amazon Mechanical Turk (MTurk) and Prolific, which now supply participants for a substantial share of published experimental work \citep{palan2018prolific,yeager2011comparing}. Yet the \emph{representativity} of these samples --- the degree to which a sample's joint distribution on outcome-relevant covariates matches that of the target population --- is constrained along three dimensions. First, respondent pools are geographically concentrated in the United States and a small number of developed countries, leaving populations in emerging markets inaccessible \citep{narasimhan2015marketing}. Second, pre-assembled panels exhibit systematic demographic skew: MTurk samples overrepresent younger, more educated, and more tech-savvy individuals \citep{stewart2015average}, and the fixed nature of the respondent pool makes it difficult to achieve representativity when relevant covariates go beyond the typical demographics that everyone wants \citep{peer2017beyond}. Finally, panel participants routinely become professional survey-takers --- \citet{rand2014social} report that their MTurk sample completed an average of 20 academic surveys in the preceding week. 

The most obvious other difference between these convenient online samples and gold-standard social surveys is, of course, cost. Traditional General Social Survey waves have been estimated at around \$3.00-\$6.67 per respondent per question \citep{goel2015non,norc2024gss}. Opt-in convenience panels such as Prolific, by contrast, run to roughly \$0.10 per respondent per question, an order-of-magnitude reduction that helps explain their displacement of probability sampling in much of applied experimental work. While representivity is important for the external validity of a study, in practice researchers face budget contraints that prevent them from gathering the sample they would ideally want. Costs constrain the researcher in two ways: 

\begin{enumerate}
\item Getting the \emph{variety} of people necessary for a representative sample is expensive for two reasons: (i) probabilistic sampling from a representative frame is itself expensive and (ii) people with certain characteristics may be more expensive to recruit into a survey compared to people with other characteristics. 
\item Getting the \emph{number} of people desired for statistical power is expensive because recruitment and participation costs generally scale with each study participant. 
\end{enumerate}

It is worth noting that the first problem (representativeness) impacts the bias of the estimate while the second problem (sample size) impacts the variance of the estimate. Both lead to estimation error.

The first chapter of this thesis reframes representativity not as a binary property but as a continuous variable that can be optimised subject to budget constraints that impact both the variety of people to be recruited and the number of people to be recruited. We additionally show that by framing study recruitment as a sequential problem of simultaneous recruitment and measurement, we can dynamically allocate resources to maximize the researchers objective (reducing the error of their population estimate). The problem is convex and can be easily optimized in practice. 

Additionally, we show that by framing the problem this way we can leverage the real-time nature of the application programming interfaces (APIs) exposed by digital ad platforms and develop an algorithm that continuously estimates stratum-specific recruitment costs from observed survey completions per dollar spent and feeds these estimates into a real-time algorithm that dynamically reallocates budget toward the strata that most improves parameter estimation per unit cost. 

The approach is enabled by three features common to major digital advertising platforms: broad population coverage, precise demographic targeting, and real-time APIs that allow programmatic control of ad delivery \citep{zhang2020quota,aridor2024experiments}. 

A U.S. validation study recruits 1,500 Meta users, benchmarks their responses against the General Social Survey, the Current Population Survey, and Pew Research Center data, and achieves an overall MAD of approximately 6.1 percentage points. The chapter compares this performance against a parallel Prolific sample of 1,200 respondents (MAD of 7.1 p.p.) and LLM-generated digital twins using the Twin-2K-500 dataset \citep{toubia2025twin} (MAD of 11.1 p.p.), at a per-question-per-respondent cost of \$0.32 --- intermediate between Prolific (\$0.10) and gold-standard probability surveys such as the GSS (\$3.00+) \citep{goel2015non}. Critically, our method allows the researcher to explicitly trade off between cost and representativity, depending on their constraints and preferences. The methodology has been deployed in 33 studies across 23 countries, with some countries experiencing costs costs well under \$0.10 per question, demonstrating global flexibility.

This adaptive stratified recruitment methodology provides the foundational infrastructure for the thesis's subsequent experimental chapters. The recruitment engine enables the two complementary designs at the heart of Chapter~\ref{ch:malaria}: it produces the representative sample that underlies the cluster-randomised trial of a malaria-prevention campaign in India, and the same respondents are then re-contacted via Facebook's remarketing tools for an ecologically valid feed experiment in which treatment is delivered as real ads in participants' actual news feeds. Chapter~\ref{ch:bebbo} similarly uses Meta-based recruitment to sample the general population of caregivers in two countries rather than self-selected app-store downloaders, preserving representativity in an encouragement design that tests a real UNICEF parenting app under natural adoption conditions. In both cases, the recruitment methodology from the first chapter makes it possible to address representativity and ecological validity simultaneously: it actively optimises the sample's demographic composition for the population of interest, and it recruits them directly through the same digital infrastructure of the intervention itself. 

\section{Experiments Advertising for Public Good}
\label{sec:intro:advertising}

Chapter~\ref{ch:malaria} evaluates whether a Facebook advertising campaign promoting malaria-prevention changes actual individual behaviour in the target population in India. The chapter recruits survey respondents via Facebook ads, works directly with the advertising team of a public health non-profit to randomise the campaign by geography, and clusters respondents into treatment and control groups at the district level to estimate impacts on self-reported attitudes, behaviours, and medical outcomes. Importantly, in this chapter we introduce a novel natural field experiment design for public advertising campaigns on digital platforms, to maximize the ecological validity of impact evaluations of such campaigns beyond what was previously possible. 

This builds on the public health literature which shows that advertising campaigns in low- and middle-income countries can shape reproductive health behaviours \citep{Agha2012}, attitudes towards smoking regulation \citep{Alday2010}, and cancer-screening behaviours \citep{Wichachai2016}, though the evidence base for social-media-delivered campaigns remains thin \citep{Schmidtke2021}. In high-income settings, digital and social media campaigns have demonstrated effects on substance-abuse behaviours \citep{Evans2020}, and traditional media campaigns delivered through cluster-randomised trials have been effective for outcomes such as childhood obesity prevention \citep{Croker2012}. The chapter reports the first evidence, to its knowledge, of a digital ad campaign's effect on malaria-related outcomes. 

Social marketing, despite its stated focus on behavioural change \citep{andreasen1994social}, has in practice relied overwhelmingly on engagement metrics as proxies for impact in evaluations of social media campaigns \citep{Shawky2019}, as opposed to actual experimentation through randomization to measure causal impacts on the outcomes of interest. 

The chapter argues that while engagement can be a means to behaviour change, the end itself --- actual behaviour change --- requires direct measurement. Researchers who have sought to measure behaviour directly have done so through artefactual field experiment designs that compromise the naturalism of the treatment delivery: embedding treatment stimuli within survey platforms \citep{Henry2020,Donati2020}, or purchasing respondents from market-research panels and delivering interventions in an artificial context \citep{Evans2021}. The chapter characterises these approaches as sharing a common disadvantage: the intervention is delivered to users during their ``active and conscious participation in a study.'' In the terminology of \citet{Schmuckler2001}, the setting dimension of ecological validity is violated --- the treatment is not encountered in the natural course of the participant's daily digital life but within an explicitly surveilled research context.

Finally, it's important to consider what representativity means in the context of Malaria prevention. Malaria disproportionately affects poor, rural populations in India \citep{Dev2004}. To estimate the impact on the overall population through a change in behavior of those at-risk for Malaria, one needs the population to be representative across the covariates that matter for Malaria exposure. In our study, we find in pilot data that household dwelling type---whether a home is built of solid (durable) or non-solid materials---is a strong and easily measured correlate of malaria risk, and a practical proxy for the broader socioeconomic conditions that drive it. Additionally, we find that individuals living in the higher-risk, non-solid dwellings were systematically under-represented in naive Facebook ad targeting relative to those in solid dwellings.

Chapter~\ref{ch:malaria} leverages the recruitment algorithm from Chapter~\ref{ch:survey-sampling} to improve representativity along this proxy, using dwelling type as the operational stratification variable. It introduces a novel methodology of using ``Lookalike Audiences'' on the Meta ad platform to target users in the harder-to-reach dwelling type, despite it not being a traditional demographic available for targeting.

The chapter strongly addresses the gap in ecological validity from the literature for public good ad campaigns and digital health intervention with a novel methodology that leverages the ad platforms retargeting capabilities to create a natural field experiment design, which we term a \emph{Feed Experiment}. 

The Feed Experiment uses Facebook's remarketing tools to individually randomise treatment among previously recruited survey respondents. Treated respondents are placed into a Custom Audience and targeted with a two-week ad campaign from an entity that has no apparent connection to the surveying entity. The ads appear in participants' real Facebook news feeds, bidding in real ad auctions alongside organic content and commercial advertisements. The outcome --- whether the respondent slept under a mosquito net the previous night --- is collected via a chatbot survey sent from the unconnected surveying entity. This design restores ecological validity on all three dimensions of \citet{Schmuckler2001}: the setting is the user's actual social-media feed, the stimulus is a real ad delivered through real platform infrastructure, and the response is a behavioural self-report collected through a channel unconnected to the advertiser. 

The chapter additionally leverages a cluster-level design to directly measure the scale-up itself through another natural field experiment, which enables estimation of both direct and spillover effects on the relevant at-risk subpopulations. Together, the two designs decompose the mechanism behind any observed campaign effect: the cluster trial measures whether an at-scale campaign moves behaviour in the population it reaches, while the feed experiment tests whether the ad content itself, conditional on exposure, is persuasive across the full spectrum of the at-risk population. Critically, the results from the scale-up experiment validate the results from the feed experiment, which in turn validates the predictive power for scale-up policies that an externally valid experiment can represent. 

\section{Experiments of Digital Apps for Public Health}
\label{sec:intro:digital-apps}

Chapter~\ref{ch:bebbo} presents a framed field experiment in a context where the natural field experiment designs introduced in Chapter~\ref{ch:malaria} were cost-prohibitive. The study evaluates a parenting app, named Bebbo, for its impact on parenting knowledge, attitude, and beliefs across a study sample of 1900 parents in Serbia and Bulgaria. 
 
Early childhood development is critical for economic development, and effective intervention yields well-documented long-run gains \citep{Black2018,Gertler2015,Attanasio2022} --- yet scalable solutions remain scarce, and the evidence base for mobile parenting apps is thin \citep{Dewitt2022}. In the field-experiment taxonomy of \citet{Harrison2004}, that evidence sits toward the controlled end. Most technology-assisted parenting trials are \emph{artefactual} field experiments: a real subject pool, but a researcher-built intervention as the commodity and a researcher-structured course of use as the task \citep{Hall2015}. The closest analogue to the present study, \citet{Outhwaite2023}, goes a step further --- a \emph{framed} field experiment evaluating a real, commercially-available app, recruited online, with a parental self-efficacy outcome --- though its task is still researcher-prompted (instructed frequency, reminder emails) and its usage self-reported, on a small UK sample (N=79) with no follow-up beyond the four-week intervention.

The Bebbo study keeps the same framed design but pushes it toward the field on the two factors where prior work is most controlled: the commodity and the task. Bebbo is a free, UNICEF-developed app, publicly available across 11 countries in the Europe and Central Asia region --- a real commodity rather than a research instrument. The evaluation uses a randomized encouragement design \citep{Moayyedi2014}: treatment participants are offered the app, control participants a link to a static informational website. Participants know they are in a study --- this is what keeps the design framed rather than natural --- but the task is left entirely uncontrolled: they use the app on their own phones, on their own time, with no prompts, reminders, or engagement-contingent payments. This restores ecological validity on the usage dimension; the study does not claim to capture how users discover the app in the wild, and who chooses to engage may be selective, but it does measure how they use it once they have it --- a claim the chapter validates by benchmarking in-study engagement against analytics from users outside the study. Usage is observed objectively through app logs rather than self-report, on a sample of 1{,}900 parents across Serbia and Bulgaria.

The evaluation finds no significant effect on any of the eight outcomes --- spanning parental knowledge, attitudes, and practices --- under either intent-to-treat or treatment-on-the-treated analysis. The engagement data explain why: 76\% of compliers abandoned the app after the first day and only 3\% used it on more than three days, and external analytics show even lower retention (86\% abandoning after one day), so the null reflects the app, not the design. This is where the usage-dimension ecological validity pays off: the null reveals what a real parenting app, as currently designed, does in the hands of real parents.

Two further findings speak to the thesis's other contributions in addition to the null findings. The first is methodological: knowledge scores rose in both arms between baseline and endline \citep{Stantcheva2023} --- the pre-post instrument itself primed respondents, and this priming can overwhelm the power gain a pre-post design is meant to provide, a caution for survey-based evaluations of this kind. The second concerns representativity. Most respondents already scored well at baseline, so within the population actually sampled --- parents reachable through online ads in these two countries --- there is little room for a generic app to improve practice. Unlike the malaria study of Chapter~\ref{ch:malaria}, which defined an at-risk target population a priori and stratified recruitment toward it, the Bebbo study did neither; budget constraints led it to recruit the broad, conveniently-reachable online population instead. A sharper definition of who most needs the app --- and the stratified recruitment of Chapter~\ref{ch:survey-sampling} to reach them --- might have avoided these ceiling effects. The study thus underscores the value of the methods developed in the preceding chapters even where it does not deploy them due to budget constraints. 



\printbibliography[heading=subbibliography, title={References for Chapter \thechapter}]
